\documentclass{beamer}
\usepackage[utf8]{inputenc}
\usepackage[T1]{fontenc}
\usepackage[portuguese]{babel}
\usepackage{amsmath}
\usepackage{tikz}
% A linha abaixo foi corrigida com "decorations.pathreplacing"
\usetikzlibrary{shapes.geometric, arrows.meta, 3d, calc, decorations.pathreplacing}

% Configuração do Tema
\usetheme{Madrid}
\usecolortheme{default}

% Informações do Título
\title{Aula 3: Teoria e Prática}
\subtitle{Vetores em dues dimensoes}
\date{28 Fevereiro 2026}

\begin{document}
	
	\frame{\titlepage}
	
	% PAGE 1
	\begin{frame}{Um vetor bidimensional}
		\small
		Um vetor bidimensional é uma grandeza vetorial (ou seja, tem módulo, direção e orientação) que "vive" no plano bidimensional. Então é um vetor como no caso mono-dimensional, mas as possíveis direções são infinitas:
		\begin{columns}[T]
			% Coluna Esquerda: Vetor 1-d
			\begin{column}{0.4\textwidth}
				\centering
				\textcolor{red}{Vetor 1-d} \\
				\vspace{0.1cm}
				\begin{tikzpicture}[scale=0.8]
					\draw[thick] (0,0) -- (3,1.5);
					\draw[thick, red, ->] (0,0) -- (2,1);
					\draw[red, <-] (2,0.3) to[out=270,in=90] (2.2, -0.3) node[below, black, align=center] {\scriptsize Uma direção \\ \scriptsize possível};
				\end{tikzpicture}
			\end{column}
			
			% Coluna Central: VS
			\begin{column}{0.1\textwidth}
				\centering
				\vspace{1cm}
				\textcolor{red}{\Large VS}
			\end{column}
			
			% Coluna Direita: Vetor 2-d
			\begin{column}{0.5\textwidth}
				\centering
				\textcolor{red}{Vetor 2-d} \\
				\vspace{0.1cm}
				\begin{tikzpicture}[scale=0.55]
					% Linhas pretas (representando as direções infinitas)
					\draw[thick] (-2.5,0) -- (2.5,0);          % Horizontal
					\draw[thick] (0,-2.5) -- (0,2.5);          % Vertical
					\draw[thick] (-2,-2) -- (2,2);             % Diagonal 1
					\draw[thick] (-2,1.5) -- (2,-1.5);         % Diagonal 2
					\draw[thick] (-1.2,2.2) -- (1.5,-2.8);     % Diagonal 3 (mais inclinada)
					
					% Circulo
					\draw[yellow!80!orange, thick] (0,0) circle (1.2cm);
					
					% Vetores vermelhos sobre as retas
					\draw[thick, red, ->] (0,0) -- (0,1.5);             % Cima
					\draw[thick, red, ->] (0,0) -- (-1,0);              % Esquerda
					\draw[thick, red, ->] (0,0) -- (1.8,1.8);           % Cima-direita
					\draw[thick, red, ->] (0,0) -- (-0.8,-0.8);         % Baixo-esquerda
					\draw[thick, red, ->] (0,0) -- (-1.6,1.2);          % Cima-esquerda
					\draw[thick, red, ->] (0,0) -- (1,-1.86);           % Baixo-direita
					
					% Apontador de texto
					\draw[red, <-] (2.3,-1.2) to[out=0,in=180] (2.8, -0.5) node[right, black, align=left] {\scriptsize Infinidade \\ \scriptsize de direções \\ \scriptsize possíveis};
				\end{tikzpicture}
			\end{column}
		\end{columns}
		
		\vspace{0.2cm}
		Definimos um vetor 1-d como:
		$\overline{v} = v_x \overline{i}$ \quad \parbox[t]{0.75\textwidth}{onde $|v_x|$ é a magnitude de $\overline{v}$ e  $\overline{i}$ o vetor unitário ($|\overline{i}|=1$) que indica a direção  positiva da reta onde $\overline{v}$ "vive". $v_x$ é também a componente 'x' de $\overline{v}$ no eixo x.}
		
		\vspace{0.2cm}
		\begin{columns}[c]
			\begin{column}{0.65\textwidth}
				A posição da origem '0' é arbitrária. Defino ela como 'início' do vetor $\overline{v}$.
			\end{column}
			\begin{column}{0.35\textwidth}
				\begin{tikzpicture}[scale=0.9]
					\draw[thick, ->] (0,0) -- (3,0) node[right] {$x$};
					\filldraw (1,0) circle (1pt) node[below] {$0$};
					\draw[thick, red, ->] (1,0) -- (2.5,0) node[midway, below] {\scriptsize $v_x$};
				\end{tikzpicture}
			\end{column}
		\end{columns}
	\end{frame}
	
	
	% PAGE 2
	\begin{frame}{Vetores Bidimensionais}
		\scriptsize
		\textcolor{red}{Questão:} como definimos os vetores em 2-d?
		
		\vspace{0.1cm}
		Agora $\overline{v}$ é definido de duas componente $v_x$ e $v_y$. $v_x$ é a componente no eixo x e $v_y$ no eixo y:
		
		\begin{columns}[T]
			\begin{column}{0.2\textwidth}
				$\overline{v} = v_x \overline{i} + v_y \overline{j}$
			\end{column}
			\begin{column}{0.35\textwidth}
				\begin{tikzpicture}[scale=0.75]
					% Eixos
					\draw[thick, ->] (-0.5,0) -- (2.5,0) node[right] {$x$};
					\draw[thick, ->] (0,-0.5) -- (0,2.5) node[above] {$y$};
					% Vetor v
					\draw[thick, ->] (0,0) -- (2,1.5) node[near end, above left] {$\overline{v}$};
					% Projeções
					\draw[dashed] (2,0) node[below] {$v_x$} -- (2,1.5) -- (0,1.5) node[left] {$v_y$};
					% Unitários
					\draw[thick, red, ->] (0,0) -- (0.8,0) node[below] {$\overline{i}$};
					\draw[thick, red, ->] (0,0) -- (0,0.8) node[left] {$\overline{j}$};
					% Ângulos (raios diferentes para as setas não colidirem no vetor v)
					\draw[->] (0.6,0) arc (0:36.8:0.6) node[midway, right] {$\theta$};
					\draw[->] (0,0.7) arc (90:36.8:0.7) node[midway, above] {$\varphi$};
				\end{tikzpicture}
			\end{column}
			\begin{column}{0.45\textwidth}
				$0 \le \theta \le \frac{\pi}{2} \quad , \quad 0 \le \varphi = \frac{\pi}{2} - \theta \le \frac{\pi}{2}$
				
				\vspace{0.2cm}
				$\overline{i}$ e $\overline{j}$ são os vetores unitários ($|\overline{i}|=| \overline{j}|=1$)\\
				que definem os sentidos positivos eixos 'x' e 'y'.
			\end{column}
		\end{columns}
		
		\vspace{0.1cm}
		Indicamos os vetores também com a notação: $\overline{v} = (v_x, v_y)$
		
		\vspace{0.1cm}
		O tamanho de $\overline{v}$, indicada com $v$, 	
		representa a sua magnitude calculada com o modelo:\\
		$v = |\overline{v}| = (v_x^2 + v_y^2)^{1/2} = \sqrt{v_x^2 + v_y^2}$.
		
		\vspace{0.2cm}
		\textcolor{red}{Observação}
		
		\begin{columns}[c]
			\begin{column}{0.65\textwidth}
				O triângulo aqui é retângulo $\Rightarrow v_x = v \cos \theta \quad , \quad v_y = v \sin \theta$\\
				\vspace{0.2cm}
				Também o triângulo ``cabeça para baixo'' é retângulo $\Rightarrow v_x = v \sin \varphi \quad , \quad v_y = v \cos \varphi$
			\end{column}
			\begin{column}{0.35\textwidth}
				\begin{tikzpicture}[scale=0.7]
					% Triângulo Base e Retângulo Tracejado
					\draw[dashed] (0,0) rectangle (2.5,1.5);
					\draw[thick, ->] (0,0) -- (2.5,1.5) node[midway, below right] {$\overline{v}$};
					\node[below] at (2.5,0) {$v_x$};
					\node[left] at (0,1.5) {$v_y$};
					
					% Ângulos
					\draw (0.8,0) arc (0:31:0.8) node[midway, right] {$\theta$};
					\draw (0,0.8) arc (90:31:0.8) node[midway, above] {$\varphi$};
					
					% Marcações de ângulo reto
					\draw (2.3,0) -- (2.3,0.2) -- (2.5,0.2);
					\draw (0,1.3) -- (0.2,1.3) -- (0.2,1.5);
					
					% Setas Curvas Coloridas (cruzam da esquerda para apontar para o triângulo correto)
					\draw[yellow!80!orange, thick, ->] (-1.2, 1.2) to[out=0,in=160] (1.2, 0.3); % Seta amarela para triângulo de baixo
					\draw[green!60!black, thick, ->] (-1.2, 0.3) to[out=0,in=180] (1, 1.2); % Seta verde para triângulo complementar de cima
				\end{tikzpicture}
			\end{column}
		\end{columns}
	\end{frame}
	
	
% PAGE 3
\begin{frame}{Calculo dos cateto com seno e coseno 2}
	\scriptsize
	Já vimos que em um triangulo retangulo
	
	\begin{columns}[c]
		\begin{column}{0.75\textwidth}
			$a = c \cos \theta \quad , \quad b = c \sin \theta \quad , \quad \frac{b}{a} = \frac{c \sin \theta}{c \cos \theta} = \frac{\sin \theta}{\cos \theta} = tg \theta$
			
			\vspace{0.3cm}
			$\leadsto tg \theta = \frac{b}{a} \leadsto \theta = arctg \frac{b}{a} \quad , \quad arctg \text{ é a função inversa } tg^{-1} \text{ de } tg$
		\end{column}
		\begin{column}{0.3\textwidth}
			\begin{tikzpicture}[scale=0.9]
				\draw[thick] (0,0) -- (3,0) node[midway, below] {$a$};
				\draw[thick] (3,0) -- (3,1.2) node[midway, right] {$b$};
				\draw[thick] (0,0) -- (3,1.2);
				\draw (0.8,0) arc (0:21.8:0.8) node[midway, right] {$\theta$};
			\end{tikzpicture}
		\end{column}
	\end{columns}
	
	\vspace{0.4cm}
	Se definimos um sistema cartesiano orientado (com direções positivas e negativas) e temos a hipotenusa 'v' representar o modulo de um vetor $\overline{v}$, no calculo das componentes $v_x$ e $v_y$, é necessario considerar se as componentes são positiva o negativa.
	
	\vspace{0.2cm}
	\begin{columns}[T]
		\begin{column}{0.45\textwidth}
			\vspace{0.5cm}
			\begin{itemize}
				\item \underline{positivas}: quando estão na parte positiva do eixo;
				\item \underline{negativas}: quando estão na parte negativa do eixo;
			\end{itemize}
		\end{column}
		\begin{column}{0.55\textwidth}
			\begin{tikzpicture}[scale=0.85]
				% Eixos
				\draw[thick, ->] (-2.5,0) -- (2.5,0) node[below] {$x$};
				\draw[thick, ->] (0,-1.5) -- (0,2) node[left] {$y$};
				
				% Vetor e linhas do triângulo
				\draw[thick, ->] (0,0) -- (-2,1.5) node[midway, above right] {$v$};
				\draw[thick] (-2,0) -- (-2,1.5);
				\draw[thick] (-2,0) -- (0,0);
				
				% Chaves (Braces) com a configuração feita linha a linha
				\draw[decorate, decoration={brace, amplitude=5pt, mirror}] (-2,-0.1) -- (0,-0.1) node[midway,below=6pt] {$v_x$};
				\draw[decorate, decoration={brace, amplitude=5pt}] (-2.1,0) -- (-2.1,1.5) node[midway,left=6pt] {$v_y$};
				
				% Ângulos
				\draw[->] (0.5,0) arc (0:143.1:0.5) node[near end, above right] {$\theta$};
				\draw[->] (-0.8,0) arc (180:143.1:0.8) node[midway, right] {$\varphi$};
				
				% Textos e setas a vermelho
				\node[text width=3.5cm, align=left, font=\scriptsize] (t1) at (2.5, -0.6) {$\theta$ é o angulo \\ antihorario a partir \\ do eixo 'x'.};
				\draw[red, ->] (t1.west) to[out=180,in=0] (0.6, 0.4);
				
				\node[text width=3.5cm, align=left, font=\scriptsize] (t2) at (2.5, -2) {$\varphi$ é o angulo horario \\ a partir do eixo 'x'.};
				\draw[red, ->] (t2.west) to[out=180,in=270] (-0.5,-1) to[out=90,in=290] (-0.8,0.1);
			\end{tikzpicture}
		\end{column}
	\end{columns}
\end{frame}
	
	
% PAGE 4
\begin{frame}{Componentes e Ângulos}
	\small
	componentes $v_x$ e $v_y$, com $\theta$ e $\varphi$ como na figura anterior, \\
	são calculadas como:
	
	\vspace{0.4cm}
	$v_x = v \cos \theta = v(-\cos \varphi)$, \quad `-' porque a componente $v_x$ é negativa
	
	\vspace{0.4cm}
	$v_y = v \text{sen} \theta = v \text{sen} \varphi$, \quad \ \ `+' porque $v_y$ é componente positiva
	
	\vspace{0.8cm}
	'$\cos \theta$' é uma função par: $\cos \theta = \cos(-\theta)$, '$\text{sen} \theta$' é impar: $\text{sen} \theta = -\text{sen}(-\theta)$
	
	\vspace{0.3cm}
	\begin{tikzpicture}[scale=0.9]
		% Eixos
		\draw[thick, ->] (-1.5,0) -- (2.5,0) node[right] {$x$};
		\draw[thick, ->] (0,-1.5) -- (0,2) node[left] {$y$};
		
		% Retas
		\draw[thick] (0,0) -- (2, 1.2);
		\draw[thick] (0,0) -- (2, -1.2);
		
		% Ângulos
		\draw[->] (0.8,0) arc (0:31:0.8) node[midway, right] {$\theta$};
		\draw[->] (0.8,0) arc (0:-31:0.8) node[midway, right] {$-\theta$};
	\end{tikzpicture}
\end{frame}
	
% PAGE 5
\begin{frame}{Exemplos}
	\scriptsize
	
	% --- EXEMPLO 1 ---
	\begin{columns}[T]
		\begin{column}{0.05\textwidth}
			\tikz[baseline=-0.5ex]\node[circle,draw,inner sep=1pt] {1};
		\end{column}
		\begin{column}{0.35\textwidth}
			\begin{tikzpicture}[scale=0.7]
				\draw[thick, ->] (-1,0) -- (3,0);
				\draw[thick, ->] (0,-1) -- (0,2);
				\draw[thick, ->] (0,0) -- (2,1) node[midway, below right] {$\overline{v}$};
				\draw[dashed] (2,0) -- (2,1) -- (0,1);
				\draw[->] (0.8,0) arc (0:26.5:0.8) node[midway, right] {$\theta$};
			\end{tikzpicture}
		\end{column}
		\begin{column}{0.25\textwidth}
			\vspace{0.2cm}
			$\overline{v} = 4\overline{i} + 2\overline{j}$
		\end{column}
		\begin{column}{0.35\textwidth}
			A) Escrever $(v_x, v_y)$ ; \\
			\vspace{0.15cm}
			B) Calcular $v = |\overline{v}|$ ; \\
			\vspace{0.15cm}
			C) Calcular o $tg\theta$ ; \\
			\vspace{0.15cm}
			D) Calcular $\theta$
		\end{column}
	\end{columns}
	
	\vspace{0.2cm}
	A) $(v_x, v_y) = (4,2)$ \\
	\vspace{0.15cm}
	B) $v = (v_x^2 + v_y^2)^{1/2} = (16+4)^{1/2} = 4,47$ \\
	\vspace{0.15cm}
	C) $tg \theta = \frac{v_y}{v_x} = \frac{2}{4} = \frac{1}{2}$ \\
	\vspace{0.15cm}
	D) $\theta = arctg \frac{1}{2} = 26,56^\circ$
	
	\vspace{0.4cm}
	% --- EXEMPLO 2 ---
	\begin{columns}[T]
		\begin{column}{0.05\textwidth}
			\tikz[baseline=-0.5ex]\node[circle,draw,inner sep=1pt] {2};
		\end{column}
		\begin{column}{0.35\textwidth}
			\begin{tikzpicture}[scale=0.7]
				\draw[thick, ->] (-1,0) -- (2.5,0);
				\draw[thick, ->] (0,-1) -- (0,1.5);
				\draw[thick, ->] (0,0) -- (1.5,0.8) node[midway, above left] {$\overline{v}$};
				\draw[dashed] (1.5,0) -- (1.5,0.8) -- (0,0.8);
				\draw[->] (0.8,0) arc (0:30:0.8) node[midway, right] {$\theta$};
			\end{tikzpicture}
		\end{column}
		\begin{column}{0.6\textwidth}
			\vspace{0.1cm}
			$\theta = 30^\circ$, $|\overline{v}| = v = 1$ \\
			\vspace{0.3cm}
			A) Calcular $\overline{v} = v_x \overline{i} + v_y \overline{j}$
		\end{column}
	\end{columns}
	
	\vspace{0.1cm}
	\begin{columns}[T]
		\begin{column}{0.05\textwidth}
			A)
		\end{column}
		\begin{column}{0.95\textwidth}
			$v_x = v \cos \theta = 1 \cdot \cos 30^\circ = \frac{\sqrt{3}}{2} \approx 0,86$ \\
			\vspace{0.3cm}
			$v_y = v \text{sen} \theta = 1 \cdot \text{sen} 30^\circ = \frac{1}{2} = 0,5$
		\end{column}
	\end{columns}
	
\end{frame}
	
% PAGE 6
\begin{frame}{Exercícios}
	\scriptsize
	
	\begin{columns}[T]
		% Coluna 1: Dados
		\begin{column}{0.2\textwidth}
			\vspace{0.5cm}
			\tikz[baseline=-0.5ex]\node[circle,draw,inner sep=1pt, blue] {1}; 
			$|\overline{v}| = 5 = v$ \\
			\vspace{0.2cm}
			$\theta = 120^\circ$
		\end{column}
		
		% Coluna 2: Figura
		\begin{column}{0.35\textwidth}
			\begin{tikzpicture}[scale=0.7]
				% Eixos
				\draw[thick, ->] (-3,0) -- (2,0);
				\draw[thick, ->] (0,-1) -- (0,2);
				
				% Vetor v
				\draw[thick, ->] (0,0) -- (-2,1.5) node[midway, above left] {$\overline{v}$};
				
				% Projeções tracejadas
				\draw[dashed] (-2,0) -- (-2,1.5) -- (0,1.5);
				
				% Vetores Unitários (vermelho)
				\draw[thick, red, ->] (0,0) -- (0.8,0) node[below] {$\overline{i}$};
				\draw[thick, red, ->] (0,0) -- (0,0.8) node[right] {$\overline{j}$};
				
				% Ângulos
				\draw[->] (0.5,0) arc (0:143.1:0.5) node[near end, above right] {$\quad\theta$};
				\draw[->] (-0.8,0) arc (180:143.1:0.8) node[midway, left] {$\varphi$};
			\end{tikzpicture}
		\end{column}
		
		% Coluna 3: Perguntas
		\begin{column}{0.45\textwidth}
			\vspace{0.3cm}
			A) Calcular $v_x$ e $v_y$ ; \\
			\vspace{0.2cm}
			B) escrever $\overline{v}$ na forma com $v_x$ e $v_y$ \\
			\vspace{0.2cm}
			C) Qual é significado dos sinais das componentes?
		\end{column}
	\end{columns}
	
	\vspace{0.2cm}
	\textcolor{blue}{Sol.} \quad $\varphi = 180^\circ - 120^\circ = 60^\circ$
	
	\vspace{0.2cm}
	A) \hspace{0.2cm} $v_x = v \cos 120^\circ = 5 \cdot (-0,5) = -2,5$ \\
	\hspace{0.85cm} \raisebox{0.1cm}{$\big\vert$} $= v(-\cos 60^\circ)$ 
	
	\vspace{0.2cm}
	\hspace{0.6cm} $v_y = v \text{sen} 120^\circ = 5 \cdot \frac{\sqrt{3}}{2}$ \\
	\hspace{0.85cm} \raisebox{0.1cm}{$\big\vert$} $= v \text{sen} 60$
	
	\vspace{0.3cm}
	B) \quad $\overline{v} = v_x \overline{i} + v_y \overline{j} = (-2,5)\overline{i} + \left(\frac{5\sqrt{3}}{2}\right)\overline{j} = (2,5)(-\overline{i}) + \left(\frac{5\sqrt{3}}{2}\right)\overline{j}$
	
	\vspace{0.3cm}
	C) \quad O vetor $\overline{v}$ é ``para cima'' e ``para esquerda''. ``Para cima'' é representado \\
	\hspace{0.65cm} do o sinal positivo de $v_y$ e ``para esquerda'' é representado do o sinal negativo de $v_x$. \\
	\hspace{0.65cm} De fato, o vetor $\overline{v}$ é na direção $(-\overline{i})$ respeito o eixo x e $(\overline{j})$ respeito o eixo y.
	
\end{frame}


% PAGE 7
\begin{frame}{Exercício 2}
	\scriptsize
	
	\begin{columns}[T]
		\begin{column}{0.05\textwidth}
			\tikz[baseline=-0.5ex]\node[circle,draw,inner sep=1pt, blue] {2};
		\end{column}
		\begin{column}{0.25\textwidth}
			Seja $\overline{v} = (3,-4)$
		\end{column}
		\begin{column}{0.7\textwidth}
			A) Desenhar $\overline{v}$ com os eixos 'x' e 'y'; (eixos cartesianos) \\
			\vspace{0.1cm}
			B) Encontrar $\theta$ e $\varphi$?
		\end{column}
	\end{columns}
	
	\vspace{0.4cm}
	\textcolor{blue}{Sol.} A) \quad $\overline{v} = v_x \overline{i} + v_y \overline{j} = 3\overline{i} + (-4)\overline{j} = 3\overline{i} + 4(-\overline{j})$
	
	\vspace{0.3cm}
	\begin{columns}[T]
		% Coluna do Desenvolvimento (Alínea B)
		\begin{column}{0.65\textwidth}
			B) \hspace{0.1cm} $\bullet \quad \displaystyle tg \theta = \frac{v_x}{v_y} = -\frac{3}{4} \leadsto \theta = \text{arctg}\left(-\frac{3}{4}\right) = -36,87^\circ$ \\
			\vspace{0.1cm}
			\hspace{0.8cm} $\leftarrow$ '-' porque $\theta$ é horario
			
			\vspace{0.4cm}
			\hspace{0.5cm} $\bullet \quad \displaystyle tg \varphi = \frac{v_y}{v_x} = -\frac{4}{3} \leadsto \varphi = \text{arctg}\left(-\frac{4}{3}\right) = -53,13^\circ$ \\
			\vspace{0.1cm}
			\hspace{0.8cm} $\leftarrow$ '-' porque $\varphi$ é horario
			
			\vspace{0.4cm}
			Também \quad $\displaystyle \varphi = -\frac{\pi}{2} - (\theta) = -\frac{\pi}{2} - (-36,87^\circ) = -53,13^\circ$
		\end{column}
		
		% Coluna do Gráfico Principal (Alínea A)
		\begin{column}{0.35\textwidth}
			\begin{center}
				\begin{tikzpicture}[scale=0.6]
					% Eixos
					\draw[thick, ->] (-1.5,0) -- (4,0);
					\draw[thick, ->] (0,-4.5) -- (0,2);
					
					% Vetor
					\draw[thick, ->] (0,0) -- (3,-4) node[midway, below right] {$\overline{v}$};
					\draw[dashed] (3,0) -- (3,-4) -- (0,-4);
					
					% Chaves e valores
					\draw[decorate, decoration={brace, amplitude=5pt}] (0,0.1) -- (3,0.1) node[midway, above=4pt] {$v_x$};
					\node[above right] at (3,0) {$3$};
					
					\draw[decorate, decoration={brace, amplitude=5pt}] (-0.1,-4) -- (-0.1,0) node[midway, left=4pt] {$v_y$};
					\node[below left] at (0,-4) {$-4$};
					
					% Ângulos
					\draw[->] (1.5,0) arc (0:-53.1:1.5) node[midway, right] {$\varphi$};
					\draw[->] (0,-2) arc (-90:-53.1:2) node[midway, below left] {$\theta$};
				\end{tikzpicture}
			\end{center}
		\end{column}
	\end{columns}
\end{frame}



	
% PAGE 8
\begin{frame}{Soma entre vetores bi-dimensionais}
	\scriptsize
	
	Indicamos com $\overline{v}$ e $\overline{w}$ os vetores de dois deslocamentos. $\overline{v} = v_x \overline{i} + v_y \overline{j} \quad , \quad \overline{w} = w_x \overline{i} + w_y \overline{j}$
	
	\vspace{0.2cm}
	O vetor soma \quad $\overline{u} = \overline{v} + \overline{w}$ \quad é: 
	$\overline{u} = u_x \overline{i} + u_y \overline{j} = v_x \overline{i} + v_y \overline{j} + w_x \overline{i} + w_y \overline{j}$ $= (v_x + w_x)\overline{i} + (v_y + w_y)\overline{j}$
	
	\vspace{0.1cm}
	$\leadsto \quad u_x = v_x + w_x \quad , \quad u_y = v_y + w_y$
	
	\vspace{0.3cm}
	\begin{columns}[c]
		\begin{column}{0.25\textwidth}
			Graficamente
			\begin{center}
				\begin{tikzpicture}[scale=0.55]
					\draw[thick, ->] (0,0) -- (1.5,1) node[midway, above left] {$\overline{v}$};
					\draw[thick, ->] (1.5,1) -- (3,0) node[midway, above right] {$\overline{w}$};
					\draw[thick, red, ->] (0,0) -- (3,0) node[midway, below] {$\overline{u}$};
				\end{tikzpicture}
			\end{center}
		\end{column}
		\begin{column}{0.75\textwidth}
			\begin{block}{}
				Mais, representamos graficamente a soma $\vec{u} = \vec{v} + \vec{w}$:\\
				\tikz[baseline=-0.5ex]\node[circle,fill=blue!70!black,text=white,inner sep=1pt] {1}; Decidir a origem de $\vec{v}$ (início);\\
				\tikz[baseline=-0.5ex]\node[circle,fill=blue!70!black,text=white,inner sep=1pt] {2}; Fazer coincidir a \textbf{extremidade} de $\vec{v}$ com a \textbf{origem} de $\vec{w}$;\\
				\tikz[baseline=-0.5ex]\node[circle,fill=blue!70!black,text=white,inner sep=1pt] {3}; Para obter $\vec{u}$, ligue a \textbf{origem} de $\vec{v}$ com a \textbf{extremidade} de $\vec{w}$.
			\end{block}
		\end{column}
	\end{columns}
	
	\vspace{0.1cm}
	Vetor diferencia \quad $\overline{u} = \overline{v} - \overline{w} \ \ :$ \quad
	\begin{tikzpicture}[baseline=-0.5ex, scale=0.5]
		\draw[thick, ->] (0,0) -- (1,1) node[near end, above left] {$\overline{v}$};
	\end{tikzpicture}
	\quad
	\begin{tikzpicture}[baseline=-0.5ex, scale=0.5]
		\draw[thick, ->] (0,1) -- (1.5,0) node[near start, above right] {$\overline{w}$};
	\end{tikzpicture}
	
	\vspace{0.1cm}
	\begin{columns}[c]
		\begin{column}{0.65\textwidth}
			$-\overline{w}$ é \quad
			\begin{tikzpicture}[baseline=-0.5ex, scale=0.5]
				\draw[thick, ->] (2,0) -- (0,1) node[midway, above right] {$-\overline{w}$};
			\end{tikzpicture}
			\quad $\leadsto$ \parbox[c]{4cm}{usando a regra grafica de soma\\$\overline{u} = \overline{v} + (-\overline{w})$}
		\end{column}
		\begin{column}{0.35\textwidth}
			\begin{tikzpicture}[scale=0.55]
				% Triângulo da diferença de vetores (como na nota)
				\draw[thick, ->] (1,0) -- (2.5,1.5) node[midway, below right] {$\overline{v}$};
				\draw[thick, ->] (2.5,1.5) -- (0.5,2.5) node[midway, above right] {$-\overline{w}$};
				\draw[thick, ->] (1,0) -- (0.5,2.5) node[midway, left] {$\overline{u}$};
			\end{tikzpicture}
		\end{column}
	\end{columns}
	
\end{frame}
	
% PAGE 9
\begin{frame}{Propriedades dos vetores}
	\scriptsize
	
	\textcolor{blue}{Propriedades dos vetores}
	\vspace{0.1cm}
	\begin{itemize}
		\item[a)] $\overline{v} + \overline{w} = \overline{w} + \overline{v} \quad \text{(pr. comutativa)}$
		\item[b)] $\alpha \overline{v} \text{ , } \alpha \text{ escalar , é } \quad \alpha \overline{v} = \alpha (v_x, v_y) = (\alpha v_x, \alpha v_y)$
		\item[c)] $\overline{v} + (\overline{w} + \overline{u}) = (\overline{v} + \overline{w}) + \overline{u} \quad \text{(pr. associativa)}$
	\end{itemize}
	
	\vspace{0.2cm}
	\textcolor{blue}{Exercício:} \quad $\overline{v} = (2,0) \quad , \quad \overline{w} = (-2,1)$ .
	

	\begin{columns}[T]
		% Coluna Esquerda: Soluções e Gráficos
		\begin{column}{0.45\textwidth}
			\textcolor{blue}{Sol.} \\
			A) $(u_x, u_y) = (v_x + w_x, v_y + w_y)$ \\
			\hspace*{0.4cm} $= (2-2, 0+1) = (0,1)$
			
			\vspace{0.2cm}
			B) \hspace{0.1cm}
			\begin{tikzpicture}[baseline=(current bounding box.center), scale=0.55]
				\draw[thick, ->] (-1,0) -- (3,0);
				\draw[thick, ->] (0,-1) -- (0,2);
				\draw[thick, red, ->] (0,0) -- (2,0) node[midway, below] {$\overline{v}$};
				\draw[thick, red, ->] (2,0) -- (0,1) node[midway, above right] {$\overline{w}$};
				\draw[thick, green!60!black, ->] (0,0) -- (0,1) node[midway, left] {$\overline{u}$};
			\end{tikzpicture}
			

			\hspace*{0.4cm}
			\begin{tikzpicture}[baseline=(current bounding box.center), scale=0.55]
				\draw[thick, ->] (-3,0) -- (5,0);
				\draw[thick, ->] (0,-2) -- (0,2);
				\draw[thick, red, ->] (0,0) -- (-2,1) node[midway, above left] {$\overline{w}$};
				\draw[thick, red, ->] (0,0) -- (2,0) node[midway, above] {$\overline{v}$};
				\draw[thick, red, ->] (2,0) -- (4,-1)node[midway, above ] {-$\overline{w}$};;
				\draw[thick, green!60!black, ->] (0,0) -- (4,-1) node[midway, below right] {$\overline{u}$};
			\end{tikzpicture}
			
\vspace{0.2cm}


		\end{column}
		
		% Coluna Direita: Perguntas do Exercício
		\begin{column}{0.55\textwidth}
			A) Calcular $\overline{u} = \overline{v} + \overline{w}$ ;
			
			\vspace{0.1cm}
			B) Representar graficamente o vetor soma $\overline{u} = \overline{v} + \overline{w}$ e o vetor diferencia $\overline{u} = \overline{v} - \overline{w}$ ;
			
			\vspace{0.3cm}
			C) Calcular o modulo de $\overline{u}$.
			
			\vspace{0.3cm}
			D) Desenhar $-\overline{u}$.
			
			\vspace{0.3cm}
			E) Calcular o angulo antihorario $\theta$ e o angulo horario $\varphi$. Onde $\theta$ e $\varphi$ são os angulos associados a o vetor $\overline{u}$ e que comecam do eixo 'x'.
		\end{column}
	\end{columns}
	
\textcolor{red}{$-\overline{w} = -(-2,1) = (2,-1)$}, $\textcolor{green!60!black}{	\overline{u} = \overline{v} - \overline{w} = \overline{v} + (-\overline{w})=(2+2,-1)= (4,-1) }$
	
\vspace{0.2cm}	
\end{frame}


% PAGE 10
\begin{frame}{}
	\scriptsize
	
	C) \quad $|\overline{u}| = \sqrt{0+1} = 1 \quad , \quad |\overline{u}| = \sqrt{16+1} = \sqrt{17}$
	
	\vspace{0.3cm}
	D)
	\begin{center}
		\begin{tikzpicture}[scale=0.6]
			\draw[thick, ->] (-2.5,0) -- (2.5,0);
			\draw[thick, ->] (0,-2) -- (0,2);
			\draw[thick, green!60!black, ->] (0,0) -- (0,-1.5) node[right] {$-\overline{u}$};
		\end{tikzpicture}
		\hspace{1.5cm}
		\begin{tikzpicture}[scale=0.6]
			\draw[thick, ->] (-2.5,0) -- (2.5,0);
			\draw[thick, ->] (0,-2) -- (0,2);
			\draw[thick, green!60!black, ->] (0,0) -- (-2, 0.5) node[above right] {$-\overline{u}$};
		\end{tikzpicture}
	\end{center}
	
	\vspace{0.3cm}
	E)
	\begin{columns}[c]
		% E - Left Graph (Primeiro Vetor)
		\begin{column}{0.35\textwidth}
			\begin{center}
				\begin{tikzpicture}[scale=0.65]
					\draw[thick, ->] (-2,0) -- (2.5,0);
					\draw[thick, ->] (0,-2.5) -- (0,2.5);
					\draw[thick, green!60!black, ->] (0,0) -- (0,1.5) node[below left] {$\overline{u}$};
					
					% Arco laranja (theta) e texto posicionado manualmente
					\draw[->, orange, thick] (0.5,0) arc (0:90:0.5);
					\node[orange, right] at (0.4, 0.6) {$\theta = \frac{\pi}{2}$};
					
					% Arco amarelo (phi) e texto posicionado manualmente
					\draw[->, purple!80!orange, thick] (0.8,0) arc (0:-270:0.8);
					\node[purple!80!orange, left] at (-0.8, 1) {$-\frac{3}{2}\pi = \varphi$};
				\end{tikzpicture}
			\end{center}
		\end{column}
		
		% E - Right Graph & Equations (Segundo Vetor)
		\begin{column}{0.65\textwidth}
			\begin{columns}[c]
				\begin{column}{0.45\textwidth}
					\begin{center}
						\begin{tikzpicture}[scale=0.65]
							\draw[thick, ->] (-2,0) -- (2.5,0);
							\draw[thick, ->] (0,-2.5) -- (0,2);
							\draw[thick, green!60!black, ->] (0,0) -- (2,-0.5) node[below] {$\overline{u}$};
							
							% Arco amarelo (phi) com raio maior para não colidir
							\draw[->, purple!80!orange, thick] (1.4,0) arc (0:-14:1.4) node[right] {$\varphi$};
							
							% Arco laranja (theta) centrado à esquerda
							\draw[->, orange, thick] (0.7,0) arc (0:346:0.7) node[midway, left] {$\theta$};
						\end{tikzpicture}
					\end{center}
				\end{column}
				\begin{column}{0.55\textwidth}
					$\overline{u} = (4,-1)$ \\
					\vspace{0.2cm}
					$\displaystyle tg \varphi = -\frac{1}{4}$ \\
					\vspace{0.2cm}
					$\begin{aligned}
						\varphi &= \text{arctg}\left(-\frac{1}{4}\right) \\
						&\approx -14^\circ
					\end{aligned}$ \\
					\vspace{0.2cm}
					$\leadsto \theta \approx 360^\circ - 14^\circ = 346^\circ$
				\end{column}
			\end{columns}
		\end{column}
	\end{columns}
	
\end{frame}


% PAGE 11
\begin{frame}{Normalização de vetores}
	\scriptsize
	
	Se tenho um vetor $\overline{v} = v_x \overline{i} + v_y \overline{j}$ , como faço a encontrar o vetor unitário $\overline{e}_v$ na orientação de $\overline{v}$? Ou seja o vetor $\overline{e}_v$ $\underset{\textcolor{violet}{a)}}{\underline{\text{na mesma orientação}}}$ de $\overline{v}$ e tal que $\underset{\textcolor{violet}{b)}}{\underline{|\overline{e}_v| = 1}}$.
	
	\vspace{0.2cm}
	$\overline{e}_v = (e_{v,x} , e_{v,y}) \quad , \quad e_{v,x} = \frac{v_x}{|\overline{v}|} \quad , \quad e_{v,y} = \frac{v_y}{|\overline{v}|}$
	
	\vspace{0.3cm}
	Verificamos a suas propriedades \textcolor{violet}{a)} e \textcolor{violet}{b)} :
	
	\vspace{0.2cm}
	\textcolor{violet}{a)} \quad $\displaystyle \overline{e}_v = \left( \frac{v_x}{|\overline{v}|} , \frac{v_y}{|\overline{v}|} \right) \tikz[baseline=(eq.base)]{\node[inner sep=0pt] (eq) {$=$}; \draw[<-, shorten <=2pt] (eq.south) -- ++(0,-0.4) node[below, font=\tiny, align=center] {$|\overline{v}|$ é um escalar positivo \\ $\Rightarrow$ não muda a orientação \\ respeito a $\overline{v}$};} \frac{1}{|\overline{v}|} (v_x, v_y) = \left( \frac{1}{|\overline{v}|} \right) \overline{v}$
	
	\vspace{0.6cm}
	\begin{columns}[T]
		% Cálculos da propriedade b)
		\begin{column}{0.65\textwidth}
			\textcolor{violet}{b)} \quad $\displaystyle |\overline{e}_v| = \sqrt{e_{v,x}^2 + e_{v,y}^2} = \sqrt{\frac{v_x^2}{|\overline{v}|^2} + \frac{v_y^2}{|\overline{v}|^2}} = \sqrt{\frac{1}{|\overline{v}|^2} (v_x^2 + v_y^2)}$
			
			\vspace{0.3cm}
			\hspace{0.5cm} $\displaystyle = \frac{1}{|\overline{v}|} \sqrt{v_x^2 + v_y^2} = \frac{|\overline{v}|}{|\overline{v}|} = 1$
		\end{column}
		
		% Gráfico
		\begin{column}{0.35\textwidth}
			\vspace{-0.5cm}
			\begin{center}
				\begin{tikzpicture}[scale=0.7]
					\draw[thick, ->] (-1.5,0) -- (2.5,0);
					\draw[thick, ->] (0,-2.5) -- (0,1.5);
					\draw[thick, ->] (0,0) -- (2,-2) node[midway, below left] {$\overline{v}$};
					\draw[thick, red, ->] (0,0) -- (0.8,-0.8) node[midway, right] {$\overline{e}_v$};
				\end{tikzpicture}
			\end{center}
		\end{column}
	\end{columns}
	

	\textcolor{red}{ Chamamos $\overline{e}_v$ o vetor normalizado / vetor unitário associado a $\overline{v}$}
\end{frame}

	
% PAGE 12
\begin{frame}{Exercício Prático}
	\scriptsize
	
	\textcolor{blue}{\underline{Exercício}} \hspace{0.5cm} Seja $\overline{v} = \left(\frac{3}{2} , -\frac{1}{2}\right)$ , Calcule o vetor unitário $\overline{e}_v$.
	
	\vspace{0.6cm}
	\textcolor{blue}{Sol.} \hspace{0.5cm} $\displaystyle |\overline{v}| = \sqrt{\left(\frac{3}{2}\right)^2 + \left(-\frac{1}{2}\right)^2} = \sqrt{\frac{9}{4} + \frac{1}{4}} = \sqrt{\frac{10}{4}} = \frac{\sqrt{10}}{\sqrt{4}} = \frac{\sqrt{10}}{2}$
	
	\vspace{0.7cm}
	$\displaystyle e_{v,x} = \frac{\frac{3}{2}}{\frac{\sqrt{10}}{2}} = \frac{3}{2} \cdot \frac{2}{\sqrt{10}} = \frac{3}{\sqrt{10}} \quad , \quad e_{v,y} = \frac{-\frac{1}{2}}{\frac{\sqrt{10}}{2}} = -\frac{1}{2} \cdot \frac{2}{\sqrt{10}} = -\frac{1}{\sqrt{10}}$
	
	\vspace{0.7cm}
	$\displaystyle \overline{e}_v = e_{v,x} \overline{i} + e_{v,y} \overline{j} = \frac{3}{\sqrt{10}} \overline{i} + \left(-\frac{1}{\sqrt{10}}\right) \overline{j} = \frac{3}{\sqrt{10}} \overline{i} + \frac{1}{\sqrt{10}} (-\overline{j})$
	
	\vspace{0.8cm}
	\textcolor{blue}{\underline{Exercício para casa}:} estudar propriedade de seno e coseno, e seus valores, \\
	(Pode usar Wikipedia ou um modelo de AI com `thinking mode' $\leftarrow$ modalidade pensamento).
\end{frame}

	
\end{document}